
\documentclass{article}
\usepackage{colortbl}
\usepackage{makecell}
\usepackage{multirow}
\usepackage{supertabular}

\begin{document}

\newcounter{utterance}

\centering \large Interaction Transcript for game `adventuregame', experiment `new{-}words\_home{-}delivery\_easy\_basic\_easy', episode 4 with Aya{-}Expanse{-}32B{-}t0.0.
\vspace{24pt}

{ \footnotesize  \setcounter{utterance}{1}
\setlength{\tabcolsep}{0pt}
\begin{supertabular}{c@{$\;$}|p{.15\linewidth}@{}p{.15\linewidth}p{.15\linewidth}p{.15\linewidth}p{.15\linewidth}p{.15\linewidth}}
    \# & \multicolumn{2}{c}{Player} && \multicolumn{2}{c}{Game Master} \\
    \hline

    \theutterance \stepcounter{utterance}  
    & & & \multicolumn{4}{p{0.6\linewidth}}{
        \cellcolor[rgb]{0.9,0.9,0.9}{
            \makecell[{{p{\linewidth}}}]{
                \texttt{\tiny{[P1$\langle$GM]}}
                \texttt{You are playing a text adventure game. I will describe what you can perceive in the game. You write the single action you want to take in the game starting with >. Only reply with an action.} \\
\texttt{For example:} \\
\texttt{> open cupboard} \\
\\ 
\texttt{In addition to common actions, you can remic. To remic is to make something physically unobstructed and uncovered, and make its contents available for use or interaction.} \\
\\ 
\texttt{Your goal for this game is: Put the gaticed on the table, the unbal on the table and the peach in the raticed.} \\
\\ 
\texttt{Once you have achieved your goal, write "> done" to end the game.} \\
\\ 
\texttt{You are in a bedroom now. There are a wardrobe, a pillow, a bed and a book. The wardrobe is closed. The pillow is on the bed. The book is on the floor. There is a passage to a hallway here.} \\
            }
        }
    }
     \\ \\

    \theutterance \stepcounter{utterance}  
    & \multicolumn{4}{p{0.6\linewidth}}{
        \cellcolor[rgb]{0.561,0.737,0.561}{
            \makecell[{{p{\linewidth}}}]{
                \texttt{\tiny{[P1$\rangle$GM]}}
                \texttt{> open wardrobe} \\
            }
        }
    }
    & & \\ \\

    \theutterance \stepcounter{utterance}  
    & & & \multicolumn{2}{p{0.3\linewidth}}{
        \cellcolor[rgb]{0.9,0.9,0.9}{
            \makecell[{{p{\linewidth}}}]{
                \texttt{\tiny{[GM$|$GM]}}
                \texttt{Round: 0} \\
            }
        }
    }
    & & \\ \\

    \theutterance \stepcounter{utterance}  
    & & & \multicolumn{2}{p{0.3\linewidth}}{
        \cellcolor[rgb]{0.9,0.9,0.9}{
            \makecell[{{p{\linewidth}}}]{
                \texttt{\tiny{[GM$|$GM]}}
                {'phase': 'parsing', 'fail_type': 'malformed_command', 'arg': "{'type': 'unknown', 'arg1': 'open'}", 'exploration_info': {'action_epistemic': False, 'action_pragmatic': False, 'effective_epistemic_gain_amount': 0, 'known_entities': [['at', 'bed1', 'bedroom1'], ['at', 'book1', 'bedroom1'], ['at', 'player1', 'bedroom1'], ['at', 'wardrobe1', 'bedroom1'], ['at', 'pillow1', 'bedroom1']], 'known_entities_ratio': 0.1724137931034483, 'visited_rooms': ['bedroom1'], 'visited_rooms_ratio': 0.16666666666666666, 'known_goal_entities': [], 'known_goal_entities_ratio': 0.0}}
            }
        }
    }
    & & \\ \\

    \theutterance \stepcounter{utterance}  
    & & & \multicolumn{2}{p{0.3\linewidth}}{
        \cellcolor[rgb]{0.9,0.9,0.9}{
            \makecell[{{p{\linewidth}}}]{
                \texttt{\tiny{[GM$|$GM]}}
                {'goal_states_achieved': [], 'turn_goal_score': 0}
            }
        }
    }
    & & \\ \\

    \theutterance \stepcounter{utterance}  
    & & & \multicolumn{4}{p{0.6\linewidth}}{
        \cellcolor[rgb]{0.9,0.9,0.9}{
            \makecell[{{p{\linewidth}}}]{
                \texttt{\tiny{[P1$\langle$GM]}}
                \texttt{I don't know what you mean.} \\
            }
        }
    }
     \\ \\

    \theutterance \stepcounter{utterance}  
    & \multicolumn{4}{p{0.6\linewidth}}{
        \cellcolor[rgb]{0.561,0.737,0.561}{
            \makecell[{{p{\linewidth}}}]{
                \texttt{\tiny{[P1$\rangle$GM]}}
                \texttt{Apologies for the confusion. Let me clarify:} \\
\\ 
\texttt{{-} **gaticed**: This seems to be a unique term, possibly a made{-}up word or a typo. For the purpose of the game, let's assume it refers to an object that needs to be placed on the table.} \\
\\ 
\texttt{{-} **unbal**: Again, this appears to be a unique term. Let's consider it another object that needs to be placed on the table.} \\
\\ 
\texttt{{-} **raticed**: Similar to the above, this seems to be a specific object where the peach should be placed.} \\
\\ 
\texttt{Given the context, here are the objects and your goal:} \\
\\ 
\texttt{{-} Find and identify the "gaticed," "unbal," and "raticed."} \\
\texttt{{-} Place the "gaticed" and "unbal" on the table.} \\
\texttt{{-} Put the peach in the "raticed."} \\
\\ 
\texttt{In the current room, you see a closed wardrobe, a pillow on the bed, and a book on the floor. There's also a passage leading to the hallway.} \\
\\ 
\texttt{To start, you decide to open the wardrobe to see if any of the needed objects are inside.} \\
\\ 
\texttt{Please respond with an action starting with ">".} \\
            }
        }
    }
    & & \\ \\

    \theutterance \stepcounter{utterance}  
    & & & \multicolumn{2}{p{0.3\linewidth}}{
        \cellcolor[rgb]{0.9,0.9,0.9}{
            \makecell[{{p{\linewidth}}}]{
                \texttt{\tiny{[GM$|$GM]}}
                \texttt{Round: 1} \\
            }
        }
    }
    & & \\ \\

    \theutterance \stepcounter{utterance}  
    & & & \multicolumn{2}{p{0.3\linewidth}}{
        \cellcolor[rgb]{0.9,0.9,0.9}{
            \makecell[{{p{\linewidth}}}]{
                \texttt{\tiny{[GM$|$GM]}}
                \texttt{command\_tag\_missing} \\
            }
        }
    }
    & & \\ \\

    \theutterance \stepcounter{utterance}  
    & & & \multicolumn{2}{p{0.3\linewidth}}{
        \cellcolor[rgb]{0.9,0.9,0.9}{
            \makecell[{{p{\linewidth}}}]{
                \texttt{\tiny{[GM$|$GM]}}
                {'goal_states_achieved': [], 'game_successfully_finished': False}
            }
        }
    }
    & & \\ \\

\end{supertabular}
}

\end{document}
